\section*{Chapitre \ref{ch:m1:getting-access-equities} --- Accès aux marchés~: marchés d'actions}

\begin{solution}{exo:m1:getting-access-equities:1}
$9/12\,000 = 0.075\%$ du volume consolidé~: au-dessus de 0.06~\%, en dessous de 0.20~\%~: le
premier palier, 0.0020~\$ par action. 18\,000~\$ par jour~; 378\,000~\$ pour le mois.
\end{solution}

\begin{solution}{exo:m1:getting-access-equities:2}
$23.9 \times 10^6 \times 0.0020 = \$47\,800$; $24.0 \times 10^6 \times 0.0023 =
\$55\,200$. Les cent mille dernières actions valent 7\,400~\$ par jour, soit 7.4 cents
par action, contre une ristourne de 0.23 cent, et 155\,400~\$ sur le mois.
\end{solution}

\begin{solution}{exo:m1:getting-access-equities:3}
Passif~: $-20 + 2 + 1 = -17$ cents pour 100 actions (revenu). Agressif~: $30 + 2
+ 1 = +33$. Aller-retour~: un coût net de 16 cents pour 100 actions, que
une stratégie qui entre passivement et sort agressivement doit gagner avant de
rapporter quoi que ce soit.
\end{solution}

\begin{solution}{exo:m1:getting-access-equities:4}
Le plus élevé entre le minimum du dealer, 100\,000~\$, et le montant du \omterm{def:m1:what-a-trading-firm-does:market-maker}{teneur de marché}, $300
\times 2\,500 = \$750\,000$ (sous le plafond de 1 million de dollars)~: 750\,000~\$. Endettement
agrégé maximal~: $15 \times 2$ million $= \$30$ million. Aucune des deux ne joue en
pratique~: le capital net se calcule après décote sur les positions de la société, si bien que
un portefeuille de quelques dizaines de millions consomme les 2 millions de dollars bien avant
que l'endettement ne le fasse, et la marge propre du \omterm{def:m1:getting-access-equities:brokers}{courtier compensateur} est en général
plus stricte encore.
\end{solution}

\begin{solution}{exo:m1:getting-access-equities:5}
$10\% \times 390 = 39$ minutes. Avec «~un cours acheteur \emph{ou} un cours vendeur au meilleur cours~»,
6~\% et 7~\% qui ne coïncident jamais représentent 13~\% de la journée avec un seul côté au
meilleur cours~: qualifié selon cette lecture~; s'ils coïncident toujours, 7~\%~: non
qualifié. Si la bourse calcule plutôt la moyenne des pourcentages acheteur et vendeur,
la première société affiche 6.5~\% et échoue. La méthode de mesure figure dans le règlement de marché et
doit être lue, car elle détermine le réglage du moteur de cotation.
\end{solution}

\begin{solution}{exo:m1:getting-access-equities:6}
Position dans la file~: peu de participants laissent au repos des ordres là où ils doivent payer, si bien que la
file est courte, et un ordre y est exécuté avant les longues files des
plateformes qui versent une ristourne, au moment où l'exécution a le plus de valeur. Flux~: les ordres exécutables
émis par des routeurs sensibles aux frais partent d'abord vers la \omterm{def:m1:us-equity-market-structure:makertaker}{plateforme inversée}, car elle
les rémunère, si bien que l'ordre au repos rencontre une autre contrepartie, souvent moins
informée. Ces deux effets peuvent valoir plus de 30 cents pour 100 actions.
\end{solution}

\begin{solution}{exo:m1:getting-access-equities:7}
$24 \times (27 - 23)/(27 - 20) = 13.7$ millions d'actions par jour pour une perte de 27~;
$24 \times (25 - 23)/(25 - 20) = 9.6$ millions à 25. Juste au-dessus de chaque seuil, le remplissage
rapporte quelques dollars par jour~; juste en dessous, il les fait perdre.
\end{solution}

\begin{solution}{exo:m1:getting-access-equities:8}
(i) Le taux de 31 cents est le palier supérieur, pour les membres ajoutant 1~\% du
volume consolidé, 120 millions d'actions par jour sur 12 milliards~; à 2 millions, la société est au
tarif standard de 16. (ii) Les frais de compensation et les coûts réglementaires, environ 3, sont absents.
Recalcul~: $10 + 16 - 3 = 23$ cents pour 100 actions, 96\,600~\$ par mois, et non
172\,200~\$. Il faut ajouter que les 10 cents «~avant frais~» ont vraisemblablement été
mesurés sur des exécutions que la longue file d'une plateforme versant une ristourne ne fournirait pas.
\end{solution}

\begin{solution}{pb:m1:getting-access-equities:1}
\textbf{1.} $20/12\,000 = 0.167\%$~: premier palier, 0.0020~\$~: 840\,000~\$ par mois.
\textbf{2.} $0.20\% \times 12$ milliards $= 24$ millions~: 4 millions de plus par jour.
\textbf{3.} $24 \times 0.0023 - 4 \times 0.0027 - 20 \times 0.0020$ (en millions de
dollars) $= 55\,200 - 10\,800 - 40\,000 = +\$4\,400$ par jour.
\textbf{4.} $24 \times 4/7 = 13.7$ millions d'actions par jour.
\textbf{5.} De 24 à 30 millions, les actions ajoutées par le remplissage perdent 27 et rapportent la nouvelle
ristourne de 28~: chacune s'autofinance, si bien que le volume de seuil est nul.
Le remplissage rapporte $30 \times 0.0028 - 6 \times 0.0027 - 24 \times 0.0023 = +\$12\,600$ par
jour.
\textbf{6.} $0.8 \times 0.0028 = \$0.00224$.
\textbf{7.} Client~: $0.00224 \times 20$ millions $\times 21 = \$940\,800$. Membre~:
$840\,000 - 25\,000 = \$815\,000$. La voie client est meilleure de 125\,800~\$.
\textbf{8.} Membre, au palier de 0.25~\% qui lui est propre~: $0.0028 \times 30$ millions $\times 21 -
25\,000 = \$1\,739\,000$. Client~: 1\,411\,200~\$. L'adhésion est meilleure de
327\,800~\$.
\textbf{9.} Connectivité directe à ses propres conditions (les contrôles de risque du courtier
et ses ports sont sur le chemin), l'éligibilité aux programmes des fournisseurs de liquidité,
la maîtrise des types d'ordres et des sessions, et la confidentialité~: le courtier voit
son flux d'ordres.
\textbf{10.} Agrément, capital net, déclarations, surveillance, inspections,
le règlement de marché de la bourse et les obligations d'accès au marché qui lui sont propres.
\textbf{11.} $(24 \times 21 - 23 \times 15)/6 = 26.5$ millions par jour.
\textbf{12.} $0.20\% \times 13$ milliards $= 26$ millions~: le seuil a augmenté
de 2 millions d'actions par jour sans que la société ait rien fait. Les objectifs doivent
être fixés sur une prévision du volume consolidé et suivis quotidiennement.
\textbf{13.} Au taux marginal, y compris l'effet de palier~: une stratégie dont
le volume fait franchir un seuil à la société vaut davantage pour elle que ne le montre son propre
P\&L, et le backtest doit présenter à la fois le chiffre autonome et
celui à l'échelle de la société.
\textbf{14.} Les effets de seuil rendent le volume collant~: un membre proche d'un seuil va
concentrer ses flux d'ordres sur cette bourse pour le franchir, et un membre juste au-dessus
a intérêt à y rester. Un barème continu renoncerait à cet effet d'attraction.
\textbf{15.} Il ajoute de la liquidité affichée, ce que la bourse rémunère
mais une liquidité qui existe pour la ristourne et non pour le prix est la
première à disparaître~; elle fausse en outre les statistiques de routage. Les transactions sans
autre but économique que la ristourne, ou conclues entre parties liées, basculent
dans le wash trading, ce qui est interdit.
\textbf{16.} $0.0031 - 0.0020 = 11$ cents pour 100 actions. Elle peut coter un tick plus serré
sur certaines valeurs où la plus petite société ne le peut pas, ou maintenir la même cotation
avec une marge plus épaisse~: dans une activité où l'avantage brut atteint 10 à 20 cents pour
100 actions, cela représente l'essentiel.
\textbf{17.} Les ristournes financent la fourchette que gagnent les stratégies passives~; un plafond
abaisse à la fois la ristourne et les frais d'accès, resserre les écarts
entre types de plateformes, réduit la valeur des paliers et des économies d'échelle, et déplace
les revenus du market making fondé sur les ristournes vers des stratégies qui gagnent la fourchette
elle-même (voir au \cref{ch:m1:us-equity-market-structure} pour l'état de cette
règle).
\textbf{18.} Actions par jour et par plateforme, ajoutées et retirées~; débit des messages~;
soldes moyens longs et courts et usage des titres difficiles à emprunter~; le palier que
la société atteindrait seule~; et son autre option réaliste.
\textbf{19.} \textbf{13.7 millions d'actions par jour.}
\textbf{20.} Le volume, mesuré selon la méthode choisie par la bourse.
\end{solution}

\begin{solution}{iq:m1:getting-access-equities:1}
Le \omterm{def:m1:getting-access-equities:brokers}{courtier exécutant} est le membre dont le nom et les contrôles de risque portent les
ordres de la société jusqu'à la plateforme. Le \omterm{def:m1:getting-access-equities:brokers}{courtier compensateur} est le membre de la
chambre de compensation qui règle les transactions, détient les positions et les liquidités,
les finance et appelle les marges. Il peut s'agir de la même société~; avec un accord de
\omterm{def:m1:getting-access-equities:brokers}{give-up}, une société exécute par l'intermédiaire de plusieurs courtiers et compense auprès d'un seul.
\iqlookfor{qui porte le risque de crédit~: le \omterm{def:m1:getting-access-equities:brokers}{courtier compensateur}.}
\end{solution}

\begin{solution}{iq:m1:getting-access-equities:2}
La bourse verse une ristourne à l'ordre qui était au repos (maker) et
facture l'ordre qui a été exécuté contre lui (taker), en conservant la différence.
Un \omterm{def:m1:getting-access-equities:tier}{palier de volume} accorde une meilleure ristourne ou de meilleurs frais aux membres dont le volume mensuel
dépasse un seuil, généralement une part du volume consolidé, et l'applique à
l'ensemble du mois.
\iqlookfor{l'application rétroactive, d'où les effets de seuil.}
\end{solution}

\begin{solution}{iq:m1:getting-access-equities:3}
Pour la file d'attente. Là où l'ajout coûte de l'argent, peu d'acteurs ajoutent des ordres, si bien qu'un nouvel ordre se trouve près
du début de la file et s'exécute en premier~; et les preneurs de liquidité y sont rémunérés, si bien que les routeurs sensibles aux frais
envoient d'abord leurs ordres vers cette plateforme. Une exécution plus rapide et plus certaine face à un flux
moins informé peut valoir davantage que la différence de frais.
\iqlookfor{la priorité dans la file comme ce qui est acheté.}
\end{solution}

\begin{solution}{iq:m1:getting-access-equities:4}
Par exécution, par plateforme, à l'ajout ou au retrait, par type d'ordre, au taux du palier
que la société détient réellement, avec la grille en vigueur à la date de la transaction (les
grilles changent chaque mois). Ajouter les frais de compensation et réglementaires. Pour les études de capacité, utiliser
le palier marginal que le volume propre de la stratégie atteindrait. Jamais un forfait
«~un demi-cent par action~».
\iqlookfor{grilles à la date considérée et prise en compte des paliers.}
\end{solution}

\begin{solution}{iq:m1:getting-access-equities:5}
Commencer comme client d'un courtier qui offre un accès direct et une bonne répercussion
de son palier~: quelques jours de mise en place, aucune immatriculation. Choisir le \omterm{def:m1:getting-access-equities:brokers}{courtier
compensateur} d'abord~: méthode de calcul des marges, financement et taux d'emprunt, stabilité. Ajouter la colocalisation
et les \omterm{def:m1:us-equity-market-structure:sip}{flux directs} dès que la latence commence à compter~; prévoir le budget des licences de données,
y compris la \omterm{def:m1:getting-access-equities:nondisplay}{redevance non-display}. S'immatriculer et devenir membre lorsque les ristournes
de son propre palier, l'éligibilité aux programmes et le contrôle l'emportent sur le capital, la conformité et les
frais fixes, typiquement dès quelques dizaines de millions d'actions par jour.
\iqlookfor{compensation d'abord~; l'adhésion comme décision de volume.}
\end{solution}

\begin{solution}{iq:m1:getting-access-equities:6}
Calculer le déficit en actions, compte tenu de la dernière prévision du volume consolidé,
et la valeur du palier sur l'ensemble du mois. Puis~: router davantage de
flux passifs naturels de la société vers cette plateforme~; coter un peu plus
agressivement là où la perte espérée par action est inférieure au seuil de
la \cref{prop:m1:getting-access-equities:padding}~; avancer les flux qui seraient négociés
de toute façon. À proscrire~: échanger avec soi-même ou avec une entité liée, ou passer des ordres
sans autre but que le décompte. Fixer un arrêt~: si le déficit ne peut être comblé à
un coût acceptable d'ici le 28, cesser le gonflement artificiel du volume.
\iqlookfor{le prix d'ombre, et la frontière du wash trading.}
\end{solution}
